% SUBSTITUTE for the "Stylish Article" template (originally Vel's
% SelfArx.cls, CC-BY-NC-SA — non-commercial, so it cannot be
% redistributed here).
%
% This starter uses the upstream `lebhart` document class (CTAN, part of
% the `colorist` class series by Jinwen XU, LPPL-1.3c) — a colourful
% article class with native theorem environments and multilingual
% support. See README.md for the full disclosure.
%
% ENGINE: lebhart requires XeLaTeX or LuaLaTeX (fontspec-based); see
% engine.txt. Its default font setup (TeXGyrePagellaX, texgyreheros) is
% resolved by fontspec's own .fontspec configuration files, which point
% at real OTF filenames via kpathsea — NOT a bare fontconfig family
% lookup — so it renders correctly on this image without any font
% substitution.
%
% lebhart is not bundled: it ships in TeX Live via its own CTAN package.
\documentclass{lebhart}
\usepackage{ProjLib}

\UseLanguage{English}

% WORKAROUND: this version of colorist-fancy.sty's `keyword` environment
% references \keywordname, but projlib-author only ever defines the
% plural \keywordsname -- a genuine mismatch in the upstream package
% (TeX Live 2025), not something introduced here. Define the missing
% singular so the environment does not error out.
\providecommand{\keywordname}{Keywords}

\title{Notes on Caching Strategies for Read-Heavy Web Services}
\author{Sam Okafor}
\date{\TheDate{2026-01-15}}

\begin{document}
\maketitle

\begin{abstract}
We summarise three caching strategies commonly used in read-heavy web
services --- cache-aside, read-through, and write-behind --- and give
simple guidance for choosing between them based on consistency
requirements and write volume.
\end{abstract}
\begin{keyword}
caching, web services, cache invalidation, read-through, write-behind
\end{keyword}

\section{Introduction}

Most web services read far more often than they write. When a
database query becomes a bottleneck, an in-memory cache in front of it
is usually the cheapest fix available, but the strategy used to keep
that cache populated and correct has real consequences for both
latency and data freshness.

\section{Three strategies}

\subsection{Cache-aside}

The application checks the cache first; on a miss, it reads from the
database and populates the cache itself.

\begin{theorem}[Cache-aside invariant]\label{thm:aside}
Under cache-aside, a stale read is only possible in the window between
a database write and the corresponding cache invalidation.
\end{theorem}

\Cref{thm:aside} is the reason cache-aside implementations must treat
invalidation as part of the write path, not an afterthought.

\subsection{Read-through}

The cache itself is responsible for loading missing values from the
database, so the application only ever talks to the cache.

\subsection{Write-behind}

Writes land in the cache first and are flushed to the database
asynchronously. This gives the lowest write latency but risks losing
recent writes on a cache failure.

\section{Choosing a strategy}

\begin{itemize}
\item If reads vastly outnumber writes and brief staleness is
  acceptable, cache-aside is simplest to reason about.
\item If most access already goes through a single service layer,
  read-through centralises the loading logic.
\item Write-behind should be reserved for cases where write latency
  dominates and some data loss on failure is tolerable.
\end{itemize}

\section{Conclusion}

None of these strategies is universally correct; the right choice
follows from the read/write ratio and how much staleness or data loss
a given endpoint can tolerate.

\end{document}
